\cleardoublepage
\chapter{中期检查}

\section{项目概况}
本项目围绕“从白模自动生成与迁移材质”的核心目标展开，
致力于研究并实现一种基于已有带材质模型的参考驱动方法。
项目以三维模型或二维图片数据为基础输入，包括带有完整材质信息的参考模
型以及仅包含几何结构的白模，通过分析两者在形状、结构及语义上的关联关系，
自动推断并生成符合目标模型特征的材质贴图，从而实现对白模的材质补全与增强。

在技术实现方面，本项目重点关注材质生成与迁移算法的设计与优化。
通过结合几何特征提取、纹理映射以及可能的深度学习方法，对模型表面不
同区域进行合理的材质分配，使生成结果在视觉效果与结构一致性上达到较高质量。
同时，针对大规模数据处理与复杂模型的计算需求，对算法的效率与稳定性进行优化，以提升整体系统性能。

在系统实现层面，项目将构建一个基于 Web 的交互平台。该平台支持用户上传
三维模型文件，自动执行材质生成流程，并提供结果的在线预览与下载功能。
通过可视化界面，用户可以直观地查看材质生成效果，并进行简单交互操作，
从而降低技术使用门槛，增强系统的易用性。

总体而言，本项目融合了计算机图形学与人工智能技术，面向三维内容生成
与编辑的实际需求，力求实现从算法研究到系统应用的完整闭环，为三维模
型材质自动化处理提供一种高效且可扩展的解决方案。


\section{工作进展情况}
目前，本项目已完成数据准备与初步实验验证的关键阶段。
在数据构建方面，已从公开数据源中收集约4万个三维模型，
并对其进行统一处理。具体而言，针对每个模型完成了标准化对齐，
并从六个正交方向对模型进行渲染，生成多视角图像数据。在此基础上，
引入 aesthetic-predictor-v2-5 模型对渲染结果进行质量评估，
对图像的美观度与视觉质量进行自动打分，并据此筛选出约两万多个高质量模型，作为后续训练与实验的数据基础。

在模型训练准备方面，当前正在基于官方提供的 Dataset Preparation Toolkit 
构建适用于 Trellis2 微调的数据集。整体流程涵盖数据元信息构建、
模型下载、网格与 PBR 材质提取、体素化表示转换以及潜变量编码等多个阶段。
具体而言，首先通过构建与更新 metadata 对数据进行统一管理；
随后对原始三维资产进行标准化处理，包括 mesh 与 PBR 纹理的解析与整理；
在此基础上，将模型转换为 O-Voxel 表示，以适配后续生成模型的输入需求。
进一步地，为支持模型训练，还需对处理后的数据进行多阶段 latent 编码，
包括形状（shape）、材质（PBR）以及稀疏结构（SS）等信息的编码，
以构建多模态一致的隐空间表示。同时，通过多视角渲染生成图像条件数据，
使模型能够学习图像与三维结构之间的对应关系。通过上述流程，最终形成包含几何结构、
材质信息、体素表示及多视角图像在内的统一数据格式，从而满足 Trellis2 在多视角一致性、
结构对齐以及条件生成等方面的训练需求。

在实验进展上，已完成基于单样本的初步微调实验：
选取单个模型的四个视角数据，对 Trellis2 模型进行小规模微调，
并对其多视角生成能力进行了初步验证。实验结果表明，模型在一定程度上能够学习多视角之间的关联关系，
并生成具有基本一致性的结果，为后续大规模训练提供了可行性依据。

下一阶段的工作将重点开展大规模数据训练，利用已筛选的高质量数据集对模型进行系统性微调。
同时，将设计相应的评估方案，从生成结果的视觉质量、
多视角一致性以及与原始几何结构的匹配程度等方面对模型性能进行定量与定性分析，以进一步提升模型的实用性与稳定性。

\section{问题与建议}
在当前数据处理与模型准备过程中，主要存在以下几个方面的问题。
首先，数据处理流程较长且计算开销较大，尤其是在网格处理、
PBR 贴图解析以及体素化转换阶段，对 CPU 资源消耗较高，整体处理效率仍有优化空间。
其次，在多视角渲染与数据筛选过程中，虽然已引入 aesthetic-predictor-v2-5 进行质量评估，
但该评分机制主要侧重图像层面的美学质量，对三维结构完整性与材质一致性的衡量仍不够充分，
可能导致部分结构质量较高但视觉评分偏低的模型被过滤。

此外，在构建适用于 Trellis2 微调的数据集过程中，
不同阶段（mesh、voxel、latent、multi-view rendering）
之间的数据对齐与一致性维护较为复杂，容易在流程中引入误差，
后续仍需进一步完善自动化检查与校验机制，以提升数据可靠性。

针对以上问题，建议在后续工作中重点优化以下方向：
一是提升数据处理流程的并行化程度，以减少整体处理时间；
二是引入更综合的质量评估指标，将几何质量与渲染质量结合进行筛选；
三是加强数据管线的自动化校验与日志追踪能力，确保各阶段数据的一致性与可复现性。

\section{其他}
在系统实现的前期准备阶段，已对前端相关技术方案进行了调研与选型。
前端部分采用 Vite + React 技术栈进行开发，
其中 Vite 用于提供快速的开发构建与热更新能力，
React 用于构建组件化用户界面，以提升系统的开发效率与可维护性。

在网页端三维模型展示方面，调研了基于 WebGL 的多种实现方案，
最终选定使用 Three.js 作为核心三维渲染引擎，用于实现模型的加载、
渲染与交互展示功能，包括旋转、缩放以及多视角观察等操作，从而满足用户对三维资产可视化的需求。

通过上述技术选型，已完成前端系统的整体技术路线确定，
为后续 Web 交互平台的开发与功能实现提供了稳定的技术基础与实现框架。
